% STATUS: written 2026-11-05 (Hazewinkel 2013 pp. 4--5 and the
% bibliographical notes read; Sweedler 1980 p. 29 read on the scanned page).
\section{Rings of differential operators}\label{sec:polynomial-weyl:differential-operators}

Section~\ref{sec:polynomial-weyl:weyl} gave the Weyl algebra by
generators and relations. This section builds the same algebra by
operators. The order of an operator counts how many commutators kill it,
and the recursive characterization of Grothendieck makes that count a
definition \cite[Theorem~6.2, p.~4]{Hazewinkel2013}. The characterization
is the affine case of EGA IV, propositions 16.8.8 and 16.8.9
\cite[p.~5]{Hazewinkel2013}.

\begin{theorem}[{\cite[Theorem~6.2, p.~4]{Hazewinkel2013}}]
\label{thm:pw:diffop-order}
Let $D$ be a $k$-module endomorphism of the commutative $k$-algebra $A$.
Then $D$ is a differential operator of order $\leq n$ if and only if for
every $(n+1)$-tuple of elements $x_{0},x_{1},\dots,x_{n}\in A$ the
iterated commutator
\[
\operatorname{ad}(x_{0})\operatorname{ad}(x_{1})\cdots
\operatorname{ad}(x_{n})D \;=\; 0 .
\]
\end{theorem}

The differential operators of a commutative algebra form an algebra under
composition: the product of an operator of order $\leq n$ and an operator
of order $\leq n'$ has order $\leq n+n'$
\cite[6.4, p.~5]{Hazewinkel2013}. The filtered algebra they form is
denoted $\operatorname{Diff}(A)$ \cite[6.4, p.~5]{Hazewinkel2013}. On
$\mathbb{Q}[X,Y]$ this is exactly the na\"ive algebra of finite sums of
the monomial operators $f_{i,j}\partial_{X}^{i}\partial_{Y}^{j}$; on
$\mathbb{Z}[X,Y]$ it is larger, since the divided powers
$\frac{1}{i!}\partial_{X}^{i}$ enter \cite[6.4, p.~5]{Hazewinkel2013}.

Sweedler carried the same commutator rule to noncommutative algebras and
named the right derivation \cite[p.~29]{Sweedler1980}.

\begin{remark}[{\cite[p.~29]{Sweedler1980}}]\label{rem:pw:sweedler-right}
A \emph{right derivation} is a linear map $d\colon A\to N$ with
$d(ab)=d(a)b+d(b)a$. Right differential operators of order $n$ are defined
inductively: $D^{0}_{A}(M,N)=\operatorname{Hom}_{A}(M,N)$, and
$D^{n}_{A}(M,N)$ is the set of $f\in\operatorname{Hom}(M,N)$ with
$[f,a]\in D^{n-1}_{A}(M,N)$ for all $a\in A$, where
$[f,a](m)=f(ma)-f(m)a$ \cite[p.~29]{Sweedler1980}.
\end{remark}

The module of order-one operators splits into the multiplications and the
right derivations \cite[p.~29]{Sweedler1980}, which places the right
derivation beside the derivation of
Section~\ref{sec:basic-theory:definitions} in the family tree of
Section~\ref{sec:basic-theory:family-tree}. % claim-ok: the two cited sections carry the definitions; this sentence only names their placement

The polynomial case closes the circle: the algebra of polynomial
differential operators on affine $n$-space is the Weyl algebra $A_{n}$ of
Definition~\ref{def:pw:weyl} \cite[p.~1]{BelovKanelKontsevich2007}. The
generators of the last section are the multiplications and the partial
derivatives of this one. % claim-ok: the two cited items state both sides of the identification
On the Weyl algebra every derivation is inner by
Proposition~\ref{prop:pw:weyl-inner}. The next section varies the ground
ring instead of the operators, and builds polynomials over a derivation
itself \cite[p.~480]{Ore1933}.

Further work in the differential-operator line, held in the collection:
\noindent\emph{Representations of finite-dimensional Lie algebras by first-order differential operators Some local results in the transitive case}~\cite{LibRepresentationsOfFiniteDimensional1995}.
\noindent\emph{Derivation based differential calculi for noncommutative algebras deforming a class of three dimensional spaces}~\cite{LibDerivationBasedDifferentialCalculi2018}.
\noindent\emph{Infinite-order differential operators in the module of formal generalized functions and in a ring of formal power series}~\cite{LibInfiniteOrderDifferentialOperators2022}.

For the ring of differential operators on a smooth affine variety a finite set of generators and a finite set of defining relations are given~\cite{WlGeneratorsAndDefining2010}.

In characteristic $p$ every automorphism of the ring of differential operators on a polynomial algebra is shown to extend the Frobenius, which fixes the centre that the derivations act on~\cite{WlExtensionsOfTheFrobeniusToTheR2017}.
